% !PRINTED 133
\weberpage{133}
%<<CONT
durch $p$ theilbar ist, und ein Element, dessen Grad nicht durch $p$ theilbar ist.

Nach dem Satze II, \S.\,30 giebt es $j$ und nicht mehr von einander verschiedene conjugirte Gruppen $Q, Q', Q'', \ldots Q^{(j-1)}$ in $P$, und $j$ von einander verschiedene conjugirte Gruppen $R, R', R'', \ldots R^{(j-1)}$.

Es lässt sich hiernach der Satz beweisen:

$\gamma$) \gesperrt{Jedes Element von $P$, dessen Grad durch $p$ theilbar ist, kommt in einer und nur in einer der Gruppen $R, R', R'' \ldots R^{(j-1)}$ vor.}

Die Gruppen $Q, Q', Q'' \ldots$ sind, da ihr Grad eine Primzahl ist, cyklisch, und je zwei enthalten, weil sie von einander verschieden sind, ausser dem Einheitselemente kein gemeinsames Element. Es sei
\[
Q = 1, a, a^2, \ldots a^{p-1}.
\]
Die Gruppe $R$, die aus allen der Bedingung
\[
c^{-1}Qc = Q
\]
genügenden Elementen besteht, enthält ausser den Potenzen von $a$ kein Element vom Grade $p$; und Entsprechendes gilt für die anderen Gruppen $Q', R', Q'', R'' \ldots$

Ist nun $b$ ein Element von $P$, dessen Grad durch $p$ theilbar ist und gleich $ps$ sein mag, so dass $s$ nicht durch $p$ theilbar ist, so ist $b^s$ ein Element, dessen Grad $p$ ist, und das also in einer und nur in einer der cyklischen Gruppen $Q, Q', Q'' \ldots$ vorkommt. Ist aber $b^s = a$ ein Element von $Q$, so ist
\[
b^{-1}ab = a,
\]
d. h. $b$ kommt in der Gruppe $R$ vor, und kann in keiner der anderen Gruppen, z. B. in $R'$, vorkommen, weil sonst $b^s$ auch in $Q'$ vorkäme, was nicht möglich ist, da $Q$ und $Q'$ nur das Einheitselement gemein haben. Damit ist $\gamma$) bewiesen.

$\delta$) \gesperrt{In $U$ giebt es genau $jr_2(p-1)$ Elemente, deren Grad durch $p$ theilbar ist.}

Unter $U$ haben wir die Gesammtheit der Elemente von $P$ verstanden, deren Grad ein Theiler von $pr$ ist. Nun giebt es nach $\beta$) in $R$ genau $r_2(p-1)$ Elemente aus $U$, deren Grad durch $p$ theilbar ist, und da jede der $j$ Gruppen $R, R', R'', \ldots$ an Stelle von $R$ treten kann, so folgt $\delta$) aus $\gamma$). Um also noch zu zeigen, worauf nach dem Obigen alles ankommt, dass in $U$
%CONT>>
